\begin{theorem}[给定碰撞矩约束下的截面账本尖锐上界]\label{thm:conclusion-section-ledger-collision-moment-sharp-bound}
\leanpartial{paper\_conclusion\_section\_ledger\_collision\_moment\_sharp\_bound}{the requested Lean signature is false on the uniform law when \(n \ge 1\): taking \(\pi_i = 1/(n+1)\) forces the minimal \(q\)-moment and hence \(a=b=1/(n+1)\), contradicting the required strict inequality \(a>b\)}
沿用定理 \ref{thm:conclusion-section-ledger-kl-identity} 的记号。
令 \(q\ge 2\)，并定义归一化 \(q\)-阶碰撞矩
\[
c_q:=\sum_{x\in X}\pi(x)^q\in\left[n^{1-q},\,1\right].
\]
在约束
\[
\sum_{x\in X}\pi(x)=1,
\qquad
\sum_{x\in X}\pi(x)^q=c_q,
\qquad
\pi(x)>0
\]
下，函数
\(
\sum_x\log\pi(x)
\)
的最大值由且仅由（在置换意义下）一类两层分布实现：
存在唯一 \(a>b>0\) 使得
\[
\pi^\star_q=(a,b,\dots,b),
\qquad
a+(n-1)b=1,
\qquad
a^q+(n-1)b^q=c_q.
\]
并且对任意 \(\pi\) 满足约束，均有尖锐不等式
\[
\sum_{x\in X}\log\pi(x)\ \le\ \log a+(n-1)\log b,
\]
从而（记住 \(d(x)=N\pi(x)\)）得到截面账本的尖锐上界
\[
\log|\mathcal S(F)|
=n\log N+\sum_x\log\pi(x)
\le
n\log N+\log a+(n-1)\log b.
\]
等价地，
\[
D_{\mathrm{KL}}(\nu\|\pi)\ \ge\ D_{\mathrm{KL}}(\nu\|\pi^\star_q),
\qquad
\frac1n\sum_x\log d(x)\ \le\ \log\frac{N}{n}-D_{\mathrm{KL}}(\nu\|\pi^\star_q).
\]
\end{theorem}

\begin{proof}
考虑在正单纯形上最大化 \(f(\pi)=\sum_{i=1}^{n}\log\pi_i\)，并施加约束
\(
\sum_i\pi_i=1
\)
与
\(
\sum_i\pi_i^q=c_q
\)。
引入拉格朗日函数
\[
\mathcal L(\pi,\lambda,\mu)=\sum_{i=1}^{n}\log\pi_i+\lambda\Bigl(\sum_i\pi_i-1\Bigr)+\mu\Bigl(\sum_i\pi_i^q-c_q\Bigr).
\]
一阶必要条件给出对每个 \(i\)：
\[
\frac{\partial\mathcal L}{\partial \pi_i}
\;=\;
\frac1{\pi_i}+\lambda+\mu q\,\pi_i^{q-1}
\;=\;0.
\]
右端作为 \(\pi_i\) 的函数严格单调，故任一驻点的坐标至多取两种值。
设两值为 \(a>b>0\)，且 \(a\) 出现 \(k\) 次、\(b\) 出现 \(n-k\) 次。
\par
由一阶方程在 \(a\) 与 \(b\) 上相减，得
\[
\mu q=\frac{1/b-1/a}{a^{q-1}-b^{q-1}}>0.
\]
驻点处 Hessian 为对角阵
\[
\nabla^2_{\pi\pi}\mathcal L
=
\mathrm{diag}\!\left(-\frac1{\pi_i^2}\right)+\mu q(q-1)\,\mathrm{diag}\!\left(\pi_i^{q-2}\right).
\]
在 \(a\)-坐标上的对角元为
\[
-\frac1{a^2}+\mu q(q-1)a^{q-2}.
\]
利用均值定理，
\(
a^{q-1}-b^{q-1}\le (q-1)a^{q-2}(a-b)
\)，从而
\[
\mu q(q-1)a^{q-2}
=
(q-1)a^{q-2}\frac{1/b-1/a}{a^{q-1}-b^{q-1}}
\ge
\frac{1/b-1/a}{a-b}
=\frac{1}{ab}.
\]
因此
\[
-\frac1{a^2}+\mu q(q-1)a^{q-2}
\ge
-\frac1{a^2}+\frac1{ab}
=\frac{a-b}{a^2b}
>0.
\]
若 \(k\ge 2\)，取切向量 \(\delta\) 仅支撑在两个 \(a\)-坐标上（例如 \(\delta_1=1,\delta_2=-1\)，其余为 \(0\)），则同时满足
\(
\sum_i\delta_i=0
\)
与
\(
\sum_i q\pi_i^{q-1}\delta_i=qa^{q-1}(1-1)=0
\)，故 \(\delta\) 属于约束流形的切空间。
但此时
\[
\delta^{\mathsf T}(\nabla^2_{\pi\pi}\mathcal L)\delta
=2\Bigl(-\frac1{a^2}+\mu q(q-1)a^{q-2}\Bigr)
>0,
\]
违背极大点在切空间上的二阶必要条件（应半负定）。
故必有 \(k=1\)，即“大值层”必须唯一。
\par
于是极大解（置换意义下）必为 \((a,b,\dots,b)\) 型。
由两约束方程
\(
a+(n-1)b=1
\) 与 \(
a^q+(n-1)b^q=c_q
\)
确定 \(a,b\) 的唯一性（在区间 \(a\in(1/n,1)\) 上左端严格单调），从而得到尖锐极值与唯一性。
不等式与对 \(\log|\mathcal S(F)|\) 的结论随即推出。证毕。
\end{proof}

\begin{corollary}[\(q=2\) 的显式两层解]\label{cor:conclusion-section-ledger-collision-moment-q2-explicit}
\leanverified{paper\_conclusion\_section\_ledger\_collision\_moment\_q2\_explicit}
在定理 \ref{thm:conclusion-section-ledger-collision-moment-sharp-bound} 中取 \(q=2\)，记 \(c:=c_2=\sum_x\pi(x)^2\in[1/n,1]\)。
则极值两层分布 \(\pi^\star_2=(a,b,\dots,b)\) 的参数为
\[
a=\frac{1+\sqrt{(n-1)(nc-1)}}{n},
\qquad
b=\frac{1-\sqrt{\frac{nc-1}{n-1}}}{n}.
\]
\end{corollary}

\begin{proof}
由 \(a+(n-1)b=1\) 与 \(a^2+(n-1)b^2=c\) 解二元方程组可得所示闭式表达；取 \(a>b\) 的分支对应极大点。
证毕。
\end{proof}

\endinput
